% Phần này mô tả tích hợp solver vào giao diện game đã có của nhóm.
\section{Tích hợp thuật toán vào giao diện game}
\label{sec:ui_solver}

\subsection{Phạm vi tích hợp}

Trên nền giao diện Pygame và bộ luật chơi có sẵn của nhóm, phần tích hợp
bổ sung chế độ \textit{Solve by algorithm}. Người dùng chọn BFS, DFS, UCS
hoặc A*, chạy tìm kiếm, xem kết quả đo, phát lại lời giải và xuất CSV.
Các solver dùng chung \texttt{GameRules}; giao diện không cài đặt lại luật
di chuyển, công tắc, cầu hoặc tách khối~\cite{bloxorz_ui_source}.

Luồng xử lý được tổ chức thành ba phần:
\begin{enumerate}
    \item \texttt{FrontEnd/app.py} nhận lựa chọn, lưu trạng thái bắt đầu và
    tạo tiến trình tìm kiếm.
    \item \texttt{BackEnd/search.py} chọn lớp solver và gọi
    \texttt{Benchmark} để đo quá trình khởi tạo, tìm kiếm.
    \item Kết quả được trả về UI; \texttt{GameController} lưu chuỗi hành
    động và điều khiển replay từ đúng trạng thái đã bắt đầu tìm kiếm.
\end{enumerate}

\subsection{Chạy, dừng và phát lại lời giải}

Tìm kiếm bắt đầu từ vị trí hiện tại, bao gồm trạng thái cầu và khối con
đang hoạt động. Vì vậy, trước khi so sánh các thuật toán, cần khởi động
lại cùng một màn để bảo đảm đầu vào giống nhau.

Solver chạy trong tiến trình riêng bằng cơ chế \texttt{spawn}, giúp vòng
lặp giao diện tiếp tục xử lý sự kiện. Người dùng có thể hủy tìm kiếm;
đổi màn, chơi lại, trở về chế độ thủ công hoặc đóng cửa sổ cũng thu hồi
worker. UI có thêm thời hạn giám sát để dừng worker không phản hồi và
không sử dụng kết quả cũ sau khi trạng thái chơi thay đổi.

Khi tìm được lời giải, \textit{Replay/Pause} phát hoặc tạm dừng chuỗi hành
động, còn \textit{Step} thực hiện một hành động. Tốc độ phát lại điều chỉnh
từ 1 đến 10 hành động/giây. Replay lại từ đầu khôi phục trạng thái xuất
phát của lần tìm kiếm; thao tác di chuyển thủ công loại bỏ lời giải và
số liệu cũ. Việc phát lại lời giải tự động không ghi sao hoặc mở khóa
chiến dịch thủ công.

\subsection{Hiển thị kết quả và xuất dữ liệu}

Bảng benchmark hiển thị thời gian, đỉnh bộ nhớ Python, số node mở rộng và
số hành động/chi phí lời giải. UI phân biệt bốn trạng thái:
\texttt{solved}, \texttt{no\_solution}, \texttt{limit} và \texttt{error}.
Nếu worker không trả được bộ đếm hoặc số liệu bộ nhớ, giao diện hiển thị
không có số liệu thay vì ghi nhận một phép đo bằng 0.

Nút \textit{Export CSV} nối kết quả vào \path{output/search_results.csv},
kèm tên map, trạng thái bắt đầu và chuỗi hành động. File xuất từ UI dùng
định dạng riêng; nó chưa đồng nhất với schema benchmark hàng loạt ở
Mục~\ref{sec:benchmark}. Do đó, không ghép hai loại CSV theo vị trí cột.

\subsection{Tương thích sau khi tích hợp vào main}

Khi merge nhánh UI vào \texttt{main}, các thay đổi ở \texttt{BaseSolver},
BFS và DFS cần được phối hợp. Phiên bản tích hợp giữ API
\texttt{DFS.solve()} trả \texttt{SearchResult}, giữ
\texttt{BFS.\_search()} trả node đích và số node mở rộng; lời gọi trực tiếp
\texttt{solve()} của BFS, UCS và A* trả danh sách hành động hoặc
\texttt{None}. Wrapper \texttt{BackEnd.search.solve()} chuẩn hóa cả bốn
thuật toán thành cùng lớp \texttt{SearchResult} cho UI. DFS truyền ngoại
lệ giới hạn ra wrapper để tránh báo nhầm việc dừng tìm kiếm thành lỗi.
